%----------------------------------------------------------------------------------------
%	DOCUMENT DEFINITION
%----------------------------------------------------------------------------------------

% article class because we want to fully customize the page and not use a cv template
\documentclass[a4paper,12pt]{article}

%----------------------------------------------------------------------------------------
%	PACKAGES
%----------------------------------------------------------------------------------------
\usepackage{url}
\usepackage{parskip} 	

%other packages for formatting
\RequirePackage{color}
\RequirePackage{graphicx}
\usepackage[usenames,dvipsnames]{xcolor}
\usepackage[scale=0.9]{geometry}

%tabularx environment
\usepackage{tabularx}

%for lists within experience section
\usepackage{enumitem}

% centered version of 'X' col. type
\newcolumntype{C}{>{\centering\arraybackslash}X} 

%to prevent spillover of tabular into next pages
\usepackage{supertabular}
\usepackage{tabularx}
\newlength{\fullcollw}
\setlength{\fullcollw}{0.47\textwidth}

%custom \section
\usepackage{titlesec}				
\usepackage{multicol}
\usepackage{multirow}

\titleformat{\section}{\Large\scshape\raggedright}{}{0em}{}[\titlerule]
\titlespacing{\section}{0pt}{10pt}{10pt}

%for publications
\usepackage[style=authoryear,sorting=ynt, maxbibnames=2]{biblatex}

%Setup hyperref package, and colours for links
\usepackage[unicode, draft=false]{hyperref}
\definecolor{linkcolour}{rgb}{0,0.2,0.6}
\hypersetup{colorlinks,breaklinks,urlcolor=linkcolour,linkcolor=linkcolour}
\addbibresource{citations.bib}
\setlength\bibitemsep{1em}

%for social icons
\usepackage{fontawesome5}

%----------------------------------------------------------------------------------------
%	BEGIN DOCUMENT
%----------------------------------------------------------------------------------------
\begin{document}

% non-numbered pages
\pagestyle{empty} 

%----------------------------------------------------------------------------------------
%	TITLE
%----------------------------------------------------------------------------------------
\begin{tabularx}{\linewidth}{@{} C @{}}
\Huge{Dmitrii Khorkin} \\[7.5pt]
\href{https://github.com/dimaamega}{\raisebox{-0.05\height}\faGithub\ Github} \ $|$ \ 
\href{https://linkedin.com/in/dmitry-khorkin}{\raisebox{-0.05\height}\faLinkedin\ LinkedIn} \ $|$ \ 
\href{mailto:dmitryhorkin@gmail.com}{\raisebox{-0.05\height}\faEnvelope \ Correo} \ $|$ \ 
\href{https://t.me/dkhrkn}{\raisebox{-0.05\height}\faTelegram \ Telegram} \
\end{tabularx}

%----------------------------------------------------------------------------------------
% EXPERIENCE SECTIONS
%----------------------------------------------------------------------------------------

%Interests/ Keywords/ Summary
\section{Resumen}
Hombre, 26 años. Tengo más de 6 años de experiencia en programación en diferentes dominios.
Estoy en constante aprendizaje. Tengo una sólida formación científica y amplio conocimiento en una amplia
gama de tecnologías, incluyendo informática, aprendizaje automático y desarrollo web.
Me encanta resolver problemas complejos y explicarlos mediante visualizaciones atractivas.
Disfruto creando productos maravillosos con personas para personas

%----------------------------------------------------------------------------------------
%	Experience
%----------------------------------------------------------------------------------------
\section{Experiencia Laboral}

\begin{tabularx}{\linewidth}{ @{}l r@{} }
    \textbf{Evidently AI, Inc., Desarrollador de Software} & \hfill 1+ año: Septiembre 2023 - Presente \\[3.75pt]
    \multicolumn{2}{@{}X@{}}{
        \begin{minipage}[t]{\linewidth}
            \begin{itemize}[nosep,after=\strut, leftmargin=1em, itemsep=3pt]
                \item[--]
                Lideré el desarrollo de UI de dos proyectos relacionados: la biblioteca de código abierto Evidently (\href{https://demo.evidentlyai.com}{sitio demo}, \href{https://github.com/evidentlyai/evidently}{github}) y Evidently \href{https://www.evidentlyai.com/register}{cloud} (desde cero)
                \item[--]
                Aceleré la carga inicial de las páginas eliminando cascadas de red
                \item[--]
                Mejoré la confiabilidad y el tiempo de desarrollo de nuevas funciones implementando un proceso personalizado de sincronización entre API backend y cliente desde cero
                \item[--]
                Mejoré el proceso de pruebas creando pruebas E2E con \href{https://playwright.dev}{playwright}
                \\ \\ \underline{Stack}: \raisebox{-0.05\height}\faReact \ React, \raisebox{-0.05\height}\faCode \ MUI, \raisebox{-0.05\height}\faPython \ Python, \raisebox{-0.05\height}\faDocker \ Docker, \raisebox{-0.05\height}\faGithub \ Github Actions
            \end{itemize}
        \end{minipage}
    }
\end{tabularx}

\begin{tabularx}{\linewidth}{ @{}l r@{} }
\textbf{INPoint LLC, Ingeniero de Software} & \hfill 3.25 años: Marzo 2020 - Abril 2022, Nov 2023 - Mayo 2023 \\[3.75pt]
\multicolumn{2}{@{}X@{}}{
    \begin{minipage}[t]{\linewidth}
        \begin{itemize}[nosep,after=\strut, leftmargin=1em, itemsep=3pt]
            \item[--]
            Desarrollé el backend de \href{https://apps.apple.com/fi/app/ommie/id1510381143}{Ommie}, una red social minimalista para compartir historias personales inspiradoras.
            \href{https://ccare.stanford.edu/act-person/sasha-marchenko}{Stanford escribe} sobre nosotros
            \\ \underline{Stack}: \raisebox{-0.05\height}\faCloud \ AWS, \raisebox{-0.05\height}\faNodeJs \ Node.js, \raisebox{-0.05\height}\faDatabase \ MySQL
            \item[--]
            Ayudé a migrar el \href{https://adventures.polaris.com/}{sitio web} de Polaris Adventures de React puro (SPA) al framework Remix (SSR).
            Implementé la funcionalidad de sesión.
            Aceleré la carga inicial de las páginas (LCP y otras métricas) en aproximadamente \textbf{2.14} veces
            \\ \underline{Stack}: \raisebox{-0.05\height}\faCloud \ AWS, \raisebox{-0.05\height}\faReact \ React, \raisebox{-0.05\height}\faCode \ Remix, \raisebox{-0.05\height}\faCode \ Typescript, \raisebox{-0.05\height}\faCode \ Redis, \raisebox{-0.05\height}\faDocker \ Docker, \raisebox{-0.05\height}\faDatabase \ MySQL
            \item[--] 
            Integré el sistema CRM en los procesos de negocio existentes del banco.
            Gestioné el proceso de lanzamiento.
            Me comuniqué con múltiples equipos, incluidos desarrollo front/back, DevOps y seguridad.
            Creé un prototipo para automatizar el proceso de copia y creación de un modelo de datos complejo de CRM basado en el algoritmo DFS, lo que ahorró aproximadamente \textbf{3} días en cada implementación en un nuevo entorno
            \\ \underline{Stack}: \raisebox{-0.05\height}\faCode \ Django, \raisebox{-0.05\height}\faPython \ Python, \raisebox{-0.05\height}\faNodeJs \ Node.js, \raisebox{-0.05\height}\faEllipsisH \ Kafka, \raisebox{-0.05\height}\faDocker \ Docker, \raisebox{-0.05\height}\faDatabase \ PostgreSQL
            \item[--]
            Desarrollé el backend de la plataforma web de Etra para automatizar la selección de intercambiadores de calor, incluyendo el núcleo del sistema: algoritmos para el cálculo térmico.
            Me comuniqué con ingenieros del cliente para traducir los principios físicos que ellos comprenden en algoritmos y código.
            Desarrollé toda la infraestructura, así como sistemas personalizados de despliegue para el backend y frontend
            \\ \underline{Stack}: \raisebox{-0.05\height}\faCloud \ SelectEl, \raisebox{-0.05\height}\faNodeJs \ Node.js, \raisebox{-0.05\height}\faPython \ Python, \raisebox{-0.05\height}\faDatabase \ MySQL
        \end{itemize}
    \end{minipage}
}
\end{tabularx}

\begin{tabularx}{\linewidth}{ @{}l r@{} }
    \textbf{UNN, investigador junior} & \hfill Septiembre 2019 - Agosto 2022 \\[3.75pt]
    \multicolumn{2}{@{}X@{}}{Investigador en el \href{https://itmm.unn.ru/about/chairs/kafedra-teorii-upravleniya-i-dinamiki-sistem}{Departamento de Teoría del Control y Dinámica de Sistemas}.
    Amplié los límites del conocimiento humano sobre el proceso de sincronización en el campo de la dinámica no lineal.
    Varios de mis trabajos fueron publicados en revistas de cuartil Q1 (ver la sección 'Publicaciones' más abajo).

    Algunas presentaciones sobre mi trabajo: Licenciatura: \href{https://slides-com.translate.goog/dimaamega1/confirention-6ed465/fullscreen?\_x\_tr\_sl=auto\&\_x\_tr\_tl=en\&\_x\_tr\_hl=en\&\_x\_tr\_pto=wapp\&\_x\_tr\_hist=true}{inglés}/\href{https://slides.com/dimaamega1/confirention-6ed465/fullscreen}{ruso} tesis.
    Maestría: \href{https://slides-com.translate.goog/dmitryhorkin/master-3/fullscreen?\_x\_tr\_sl=auto\&\_x\_tr\_tl=en\&\_x\_tr\_hl=en\&\_x\_tr\_pto=wapp\&\_x\_tr\_hist=true}{inglés}/\href{https://slides.com/dmitryhorkin/master-3/fullscreen}{ruso} tesis.
    }
\end{tabularx}

%----------------------------------------------------------------------------------------
%	EDUCATION
%----------------------------------------------------------------------------------------
\section{Educación}
\begin{tabularx}{\linewidth}{@{}l X@{}}	
2020 - 2022 & \href{https://academy.yandex.com/dataschool/course/machine-learning}{Desarrollador de ML} \& parcialmente \href{https://academy.yandex.com/dataschool/course/big-data-infrastructure}{Infraestructura de Big Data} en el equivalente a un máster en \textbf{\href{https://academy.yandex.com/dataschool/}{escuela de análisis de datos de yandex}} \\
2020 - 2022 & Máster con honores en \textbf{\href{https://itmm.unn.ru/obrazovanie/magistratura/iskusstvennyj-intellekt}{UNN, informática fundamental y tecnologías de la información}} \hfill \normalsize \\
2016 - 2020 & Licenciatura en \textbf{\href{https://itmm.unn.ru/courses/pmi}{UNN, matemáticas aplicadas e informática}} \\

\end{tabularx}


%----------------------------------------------------------------------------------------
%	SKILLS
%----------------------------------------------------------------------------------------
\section{Habilidades}
\begin{tabularx}{\linewidth}{@{}l X@{}}
% Domains: & \normalsize{ \raisebox{-0.05\height} \faGlobe \ Web, \raisebox{-0.05\height} \faRobot \ ML, \raisebox{-0.05\height} \faGlasses \ Research (Nonlinear Dynamics)} \\
Lenguajes de programación: & \normalsize{\raisebox{-0.05\height}\faJs \ javascript/typescript  \raisebox{-0.05\height}\faPython \ python, \raisebox{-0.05\height}\faCode \ C/C++, \raisebox{-0.05\height}\faCode \ golang, \raisebox{-0.05\height}\faCode \ haskell (a \href{https://github.com/DimaAmega/haskell-homework/blob/main}{bit})} \\
Web: & \normalsize{ \underline{SSR}: \raisebox{-0.05\height}\faCode \ Remix, \underline{Front}: \raisebox{-0.05\height}\faCode \ html/css, \raisebox{-0.05\height}\faReact \ React, \raisebox{-0.05\height}\faCode \ Mui, \underline{Back}: \raisebox{-0.05\height}\faNodeJs \ Node.js, \raisebox{-0.05\height}\faCode \ Express } \\
Herramientas: & \normalsize{ \raisebox{-0.05\height} \faGitSquare \ Git, \raisebox{-0.05\height}\faDocker \ Docker, \raisebox{-0.05\height}\faNodeJs \ Node.js, \raisebox{-0.05\height}\faGithub \ Github Actions} \\
% Visualization: & \normalsize{ \raisebox{-0.05\height}\faJs \ Three.js, \raisebox{-0.05\height}\faJs \ PIXI.js, \raisebox{-0.05\height}\faPython \ plotly, \raisebox{-0.05\height}\faPython \ matplotlib } \\
Almacenamientos: & \normalsize{ \raisebox{-0.05\height} \faDatabase \ MySQL,  \raisebox{-0.05\height} \faDatabase \ PostgreSQL,  \raisebox{-0.05\height} \faDatabase \ MongoDB,  \raisebox{-0.05\height} \faDatabase \ S3,  \raisebox{-0.05\height} \faDatabase \ Redis} \\
Ciencias de la computación: & \normalsize{ \raisebox{-0.05\height} \faCode \ Algorithms, \raisebox{-0.05\height} \faCode \ Sistemas operativos, \raisebox{-0.05\height} \faCode \ Concurrencia \raisebox{-0.05\height} \faCode \ Programación
asíncrona, \raisebox{-0.05\height} \faCode \ Sistemas distribuidos (fundamentos)} \\

\end{tabularx}


%----------------------------------------------------------------------------------------
%	PET-PROJECTS
%----------------------------------------------------------------------------------------
\section{Proyectos personales}

\begin{tabularx}{\linewidth}{ @{}l r@{} }
    \textbf{Vectorfields} & \hfill \href{https://dimaamega.github.io/vectorfields/?x\_str=y\&y\_str=-L*y-sin(x)\&xspeed=3\&count=2000\&M\_Time\_Alive\_particle=3.95\&M\_n\_lines=24\&L=1\&skip\_welcome\&add\_scale=100}{Sitio web},  \href{https://github.com/DimaAmega/vectorfields}{Código fuente} \\[3.75pt]
    \multicolumn{2}{@{}X@{}}{Visualización de campos vectoriales 2D en el navegador}  \\
\end{tabularx}

\begin{tabularx}{\linewidth}{ @{}l r@{} }
    \textbf{Rotary States} & \hfill \href{https://github.com/unn-dynamic-systems/rotary_states}{Código fuente} \\[3.75pt]
    \multicolumn{2}{@{}X@{}}{La biblioteca minimalista para encontrar modos de rotación en sistemas EDO} \\
\end{tabularx}

\begin{tabularx}{\linewidth}{ @{}l r@{} }
    \textbf{JS Promises} & \hfill \href{https://github.com/DimaAmega/js-promise-implementation}{Código fuente} \\[3.75pt]
    \multicolumn{2}{@{}X@{}}{Implementación de "promesas" en javascript} \\
\end{tabularx}

...

Ver todos los proyectos en mi \href{https://github.com/dimaamega}{Github}

%----------------------------------------------------------------------------------------
%	PUBLICATIONS
%----------------------------------------------------------------------------------------
\section{Publicaciones}

\begin{refsection}[citations.bib]
\nocite{*}
\printbibliography[heading=none]
\end{refsection}

...

Ver todas las publicaciones en LinkedIn's \href{https://www.linkedin.com/in/dmitry-khorkin/details/publications/}{publicaciones} 
\end{document}
